\chapter{Background, Problem Statement and Objectives}

\section{Background and Rationale}\label{background-and-rationale}

Uzbekistan's health system, like those of many post-Soviet states, has
undergone a complex and protracted reform process since independence in
1991. The legacy of the Semashko model - characterized by centralized
planning, public financing and provision, and a strong hospital sector -
has shaped both the strengths and persistent challenges of the system.
Over the past three decades, the country has sought to replace
input-driven norms with outcomes-focused governance, strengthen primary
health care (PHC), and reduce out-of-pocket (OOP) payments that threaten
financial protection and equity. (Ketevan Glonti, 2013)

Despite important gains - particularly in routine immunization, maternal
and child health, and the introduction of strategic purchasing pilots -
Uzbekistan continues to face high OOP payments, gaps in financial
protection, underdeveloped quality assurance, and pronounced
rural--urban disparities. The burden of non-communicable diseases (NCDs)
has risen sharply, with ischemic heart disease, stroke, and cancers now
the leading causes of mortality. These epidemiological pressures,
combined with infrastructure and workforce constraints, have exposed the
limitations of traditional, hospital-centric models and underscored the
need for innovative, data-driven, and patient-centered approaches.
(World Health Organization, 2022)

In this context, digital health, including telemedicine platforms,
electronic health records, and e-referral systems---has emerged as a
promising instrument for health system improvement. The government's
``Concept on Health Development 2019--2025'' and subsequent strategies
have prioritized digitalization as a lever for expanding access,
improving quality, and enhancing efficiency. Early pilots in Syrdarya
region, including telemedicine clinics, e-referral systems, and digital
queue management, represent a notable inflection point in the country's
reform trajectory. (World Health Organization, 2022)

However, the digital health agenda is nascent and faces significant
barriers: legacy paper-based systems, fragmented data flows, uneven
facility readiness, limited digital literacy, and the absence of robust
governance and interoperability standards. Against this backdrop, this
thesis offers an integrated theoretical examination of health system
transformation in Uzbekistan, anchored in the specificities of the Uzbek
case and informed by comparative regional and international experience.

\section{Problem Statement and Research Questions}\label{problem-statement-and-research-questions}

\textbf{Problem Statement:}

Uzbekistan's health system faces persistent challenges in healthcare
access, service efficiency, and equity - particularly in rural areas and
among vulnerable populations. Despite ongoing reforms and digital health
pilots, gaps remain in infrastructure, coordination, and responsiveness.
Telemedicine platforms offer a promising solution, but their design must
be grounded in the lived experiences and needs of both patients and
healthcare providers. This thesis investigates how a telemedicine
platform can be developed and tested to address these challenges, based
on direct input from users and providers within the Uzbek health system.

\textbf{Research Questions:}

\begin{itemize}
\item
  What are the main barriers to healthcare access as perceived by
  patients in Uzbekistan, and how can digital tools address these
  barriers?
\item
  What are the key service delivery inefficiencies experienced by
  doctors, and what features would improve their workflow and patient
  management through a telemedicine platform?
\item
  What functional and technical requirements should a telemedicine
  platform meet to be usable, acceptable, and effective for both
  patients and providers in Uzbekistan?
\item
  How do users (patients and doctors) respond to a prototype
  telemedicine platform, and what does their feedback reveal about its
  potential to improve access and efficiency?
\end{itemize}

\section{Objectives of the Thesis}\label{objectives-of-the-thesis}

\begin{itemize}
\item
  To identify key barriers to healthcare access and service delivery
  efficiency in Uzbekistan through empirical surveys of patients and
  healthcare providers.
\item
  To translate user-reported challenges into functional and technical
  requirements for a telemedicine platform tailored to Uzbekistan's
  health system context.
\item
  To design and develop a prototype digital portal for patients and
  doctors that addresses the identified needs and supports improved
  access and workflow efficiency.
\item
  To evaluate the usability, acceptability, and perceived effectiveness
  of the telemedicine platform through pilot testing with real users.
\item
  To contribute to the theoretical and practical understanding of
  user-centered digital health design in post-Soviet health systems
  undergoing reform.
\end{itemize}

